\documentclass{beamer}
% Computer Science 1 — shared Beamer preamble.
% \input this file after \documentclass{beamer} in every Monday lecture
% deck so all decks share one visual system. Do not fork per-week copies of
% this file — edit it once here and every deck picks up the change on next
% compile. Vendored from professional_minds/presentations/beamer/theme/
% preamble.tex (same pattern semester_kickoff_week's own vendored copy
% uses), kept local rather than \input across repos so this repo stays
% self-contained and usable without a professional_minds checkout.

\usetheme{Boadilla}
\usecolortheme{seahorse}
\setbeamertemplate{navigation symbols}{}
\setbeamertemplate{footline}[frame number]

% Speaker notes: \note{...} on any frame. Compiling with the default
% `pdflatex deck.tex` produces a slides-only PDF (notes are embedded but
% hidden). Compiling with
% `pdflatex "\PassOptionsToPackage{show}{pgfpages}\input{deck.tex}"`
% or uncommenting the line below produces an instructor copy with each
% slide's notes on a second page — see README.md for both commands.
% \setbeameroption{show notes on second screen=right}

\usepackage{pgfpages}

\newcommand{\SessionTitleSlide}[4]{%
  % #1 = Week/Day (e.g. "Week 2 -- Monday"), #2 = session title,
  % #3 = essential question, #4 = primary source citation
  \begin{frame}
    \titlepage
  \end{frame}
  \begin{frame}{Essential Question}
    \Large #3

    \vspace{1em}
    \normalsize\emph{#4}
  \end{frame}
}


\title{Reading and Running AI-Generated Code}
\subtitle{CS1 Week 3 --- Wednesday, September 2, 2026}
\author{Computer Science I}
\date{}

% Screenshots are real, de-identified frames from the class recording.
\newcommand{\classshot}[2]{%
  \begin{frame}{From class: #2}
    \begin{center}
      \includegraphics[width=.98\linewidth,height=.82\textheight,keepaspectratio]{img/#1}
    \end{center}
  \end{frame}%
}

\begin{document}
\begin{frame}
  \titlepage
\end{frame}

\begin{frame}{The lesson in one sentence}
  \Large
  A chat AI tool gives you a fast first draft; it only knows what you told it. Read what it wrote, notice the gaps, and verify.
  \vfill
  \normalsize Predict. Read the code. Try one small change. Verify.
\end{frame}

\begin{frame}{Today's example: the World Bible}
  \begin{itemize}
    \item CS1 is deliberately freeform --- no textbook, four possible starter ``worlds'' to build from.
    \item Today's world: a frontier settlement with a food stock, a settler count, and a warning rule.
    \item The rule in plain English: \emph{if the food stock is below a threshold, print a warning.}
  \end{itemize}
\end{frame}

\classshot{1024_gpt_generated}{a plain-English rule turned into a first-draft Python script}

\begin{frame}{What the tool did --- and didn't --- know}
  \begin{itemize}
    \item The prompt only said: check a threshold, print a warning, use a ``settler'' voice.
    \item It did \textbf{not} say where to save the file or how to run it.
    \item That's expected: those are things you're expected to notice and supply yourself.
  \end{itemize}
\end{frame}

\begin{frame}{Reading the generated code}
  \begin{itemize}
    \item \texttt{\#} comment lines --- notes for humans, ignored by Python.
    \item \texttt{print(...)} --- sends text to the screen.
    \item \texttt{settlement\_name = "Dusty Trail"} --- a variable: a name that gives a bare value context.
  \end{itemize}
  \vspace{0.5em}
  A number alone (``8'') means nothing. \texttt{food\_stock = 8} means something.
\end{frame}

\begin{frame}[fragile]{Moving the code into a file: here-string + pipe}
\footnotesize
\begin{verbatim}
$code = @'
# generated Python code goes here
'@
$code | Set-Content -Path frontiersettlement.py -Encoding utf8
\end{verbatim}
\normalsize
\begin{itemize}
  \item \texttt{@' ... '@} is a \textbf{here-string}: a block of multi-line text.
  \item \texttt{|} is the \textbf{pipe}: what goes in one end comes out the other, unchanged.
  \item \texttt{Set-Content} writes the piped text into a file.
\end{itemize}
\end{frame}

\classshot{1744_paste_warning}{PowerShell's multi-line-paste warning --- expected, not an error}

\begin{frame}{Reading the warning}
  \begin{itemize}
    \item PowerShell is confirming you meant to run \emph{all} of the pasted text, not just the first line.
    \item Stray characters --- an extra quote, an extra space, a smart quote or emoji some tools insert --- can break a paste like this.
    \item If a pasted block looks off, reread it before choosing ``Paste anyway.''
  \end{itemize}
\end{frame}

\classshot{1815_run_output}{running \texttt{frontiersettlement.py} and reading its output}

\classshot{2045_notepad_view}{the same file, opened in Notepad --- still just text on disk}

\begin{frame}[fragile]{The \texttt{if} block, read line by line}
\footnotesize
\begin{verbatim}
if food_stock < 10:
    print("Hold onto yer hat! Our food stores are runnin' mighty low!")
\end{verbatim}
\normalsize
\begin{itemize}
  \item An \texttt{if} block runs only when its condition is true.
  \item Study this shape, then copy one line and change it to see what happens.
\end{itemize}
\end{frame}

\begin{frame}{Tool posture}
  \begin{itemize}
    \item If something is a mystery, confusing, or seems wrong: ask the tool to explain it.
    \item You stay the one deciding whether the explanation makes sense.
    \item The tool is only guessing at what will help you --- it can't read your mind.
  \end{itemize}
\end{frame}

\begin{frame}{Paired reasoning starts today}
  \begin{itemize}
    \item \textbf{Foreman:} plans on paper, doesn't type.
    \item \textbf{Worker:} has the keyboard, executes the plan.
    \item If you were a worker last week, try being the foreman this week.
    \item Keep the same partner from last week where you can.
  \end{itemize}
\end{frame}

\begin{frame}{Before Friday}
  \begin{itemize}
    \item Continue the paired work started in class.
    \item Be ready to explain what each variable represents.
    \item Be ready to explain what the \texttt{if} condition is checking.
  \end{itemize}
\end{frame}

\begin{frame}{Takeaway}
  \Huge
  \centering
  Read before you run. Ask when it's unclear. You decide what makes sense.
\end{frame}
\end{document}
